
\chapter{Estado da Arte}
\label{chap:estado_da_arte}

% Apresente o objetivo deste capítulo.
% Explique que serão analisados trabalhos acadêmicos,
% dispositivos e aplicações relacionados à proposta DoseCerta.
% Delimite os critérios utilizados para selecionar as soluções.

\section{Critérios de seleção e análise}
\label{sec:criterios_estado_arte}

% Descreva como foram selecionados os trabalhos e produtos.
% Informe as bases de pesquisa, os termos de busca,
% o período considerado e os critérios de inclusão,
% quando essas informações estiverem definidas.

\section{Dispensadores automatizados de medicamentos}
\label{sec:dispensadores}

% Apresente trabalhos e dispositivos que realizam
% armazenamento, organização ou liberação automatizada
% de medicamentos.

\subsection{Características e funcionalidades}
\label{subsec:caracteristicas_dispensadores}

% Analise funcionalidades como compartimentos,
% programação de horários, alertas e mecanismos de liberação.

\subsection{Limitações das soluções existentes}
\label{subsec:limitacoes_dispensadores}

% Discuta limitações identificadas nas fontes consultadas,
% como acessibilidade, capacidade, configuração,
% confiabilidade ou custo, quando aplicável.

\section{Aplicações digitais para organização de medicamentos}
\label{sec:aplicacoes_medicamentos}

% Analise aplicativos que permitem cadastrar medicamentos,
% consultar horários e receber lembretes.

\subsection{Funcionalidades e interação com o usuário}
\label{subsec:funcionalidades_aplicativos}

% Compare recursos e formas de interação das aplicações
% selecionadas.

\section{Soluções com recursos inteligentes}
\label{sec:solucoes_inteligentes}

% Analise soluções que incorporam recursos como
% reconhecimento de imagens, interfaces conversacionais
% ou inteligência artificial, se forem encontradas
% fontes relevantes.

\section{Comparação entre as soluções analisadas}
\label{sec:comparacao_solucoes}

% Compare as soluções com critérios comuns.
% Inclua uma tabela apenas quando houver dados
% suficientes nas fontes consultadas.

% Exemplo de estrutura para uma tabela comparativa:
%
% \begin{table}[htbp]
%     \centering
%     \caption{Comparação das soluções analisadas}
%     \label{tab:comparacao_solucoes}
%     \begin{tabular}{lccc}
%         \hline
%         Solução & Funcionalidade A & Funcionalidade B & Funcionalidade C \\
%         \hline
%         Solução 1 & ... & ... & ... \\
%         Solução 2 & ... & ... & ... \\
%         DoseCerta & ... & ... & ... \\
%         \hline
%     \end{tabular}
% \end{table}

\section{Síntese e posicionamento do DoseCerta}
\label{sec:sintese_estado_arte}

% Sintetize os principais aspectos observados.
% Explique quais necessidades ou oportunidades identificadas
% motivaram as escolhas do DoseCerta e quais características
% distinguem a proposta das soluções analisadas.
